\section{Procedimentos metodológicos}
\label{sec:procedimentos-metodologicos}

O trabalho caracteriza-se como um relato técnico de natureza aplicada e
caráter descritivo, desenvolvido no contexto de um estágio supervisionado em
desenvolvimento de software. Sua execução compreendeu o estudo da arquitetura
e das regras de negócio do sistema Banco de Especialistas, o levantamento e a
análise dos requisitos, a modelagem dos componentes necessários, a implementação
das funcionalidades, a realização de testes e a construção do ambiente de
integração contínua.

A verificação técnica envolveu testes unitários e de integração, destinados a
avaliar o comportamento das funcionalidades implementadas e a interação entre
os componentes da aplicação. Também foi verificado o funcionamento integrado
da API em ambiente de desenvolvimento conteinerizado com Docker. As atividades
foram conduzidas no contexto da arquitetura e das tecnologias já adotadas pelo
projeto, considerando suas restrições.

As atividades do módulo foram organizadas no OpenProject \cite{openproject2026},
com issues, responsáveis e prazos, o que apoiou a priorização das tarefas e o
acompanhamento do progresso ao longo do estágio.

% TODO: inserir figura do OpenProject com o backlog e as issues do módulo de solicitações

Embora o Scrum seja amplamente difundido, durante a prática do estágio
percebeu-se que o ritmo de trabalho e a natureza do projeto — desenvolvimento
focado em um módulo de um sistema legado, conduzido por um desenvolvedor com
apoio pontual de um technical lead, ambos com outras atribuições no projeto —
aproximaram-se mais dos valores e práticas do Extreme Programming (XP). O foco
não estava na execução de cerimônias rígidas de Sprints, mas nas práticas de
engenharia do XP: pequenas entregas incrementais, testes automatizados
rigorosos (Pytest), integração contínua (CI) e aplicação do padrão Additive
Change para refatoração segura do código legado. O OpenProject foi utilizado
como ferramenta de gestão do backlog e das issues, permitindo priorização e
monitoramento do fluxo de trabalho de forma adaptável.